\documentclass[a4paper,12pt]{article}
\usepackage[utf8]{inputenc}
\usepackage[T1]{fontenc}
\usepackage[ngerman]{babel}
\usepackage{amsmath}
\usepackage{amssymb}
\usepackage{enumitem}
\usepackage{geometry}
\geometry{a4paper, margin=2.5cm}
\usepackage{parskip}

\title{Einsendeaufgaben -- Aufgabe 3.2\\[0.5em]
\large Modul 61112: Lineare Algebra}
\author{}
\date{}

\begin{document}
\maketitle

\section*{Aufgabe 3.2}

\textbf{"$\Rightarrow$"-Richtung}

Sei $\beta \in \mathrm{Bil}_K(V)$ mit Rang $1$ beliebig.
Es existiert eine geordnete Basis $B = (v_1, \ldots, v_n)$.
Nach Satz 3.1.19 gilt
\[
  \beta(u,v) = k_B(u)^T M_B(\beta)\, k_B(v) \quad \forall i,j \leq n.
\]

Da $\mathrm{Rg}(\beta) = 1$, also auch $\mathrm{Rg}(M_B(\beta)) = 1$, gilt für $f: V \to V,\ v \mapsto M_B(\beta)\, k_B(v)$:
\[
  \mathrm{Rg}(f) = 1.
\]

Es existiert also eine Basis $B' = \{v_0\}$ des Unterraums $\mathrm{Bild}(f)$.
Daraus folgt:
\begin{align*}
  \beta(u,v) &= k_B(u)^T M_B(\beta)\, k_B(v) \\
             &= k_B(u)^T a_v\, k_B(v_0) \quad (\text{mit } a_v \in K) \\
             &= k_B(u)^T k_B(v_0)\, a_v = b_u \cdot a_v \quad (\text{mit } b_u \in K)
\end{align*}

Wir definieren die Linearformen
\[
  a_1: V \to K, \quad \text{und} \quad a_2: V \to K
\]
mit $a_1(v_i) := k_B(v_i)^T k_B(v_0) = b_{v_i}$
und $a_2(v_i) := k_{B'}(M_B(\beta)\, k_B(v_i)) = a_{v_i}$
für alle $v_i \in B$.

Dann gilt
\begin{align*}
  \beta(u,v) &= k_B(u)^T M_B(\beta)\, k_B(v) \\
             &= a_1(u) \cdot a_2(v).
\end{align*}

\bigskip
\textbf{"$\Leftarrow$"-Richtung}

Es existiert eine Basis $B = (v_1, \ldots, v_n)$ von $V$.
Wegen $\beta(v_i, v_j) = a_1(v_i)\, a_2(v_j)\ \forall v_i, v_j \in B$ folgt:
\begin{align*}
  M_B(\beta) &= (\beta(v_i, v_j))_{i,j} = (a_1(v_i)\, a_2(v_j))_{i,j} \\
             &= \sum_{j=1}^{n} a_2(v_j)\, (a_1(v_i))_i.
\end{align*}

Demnach hat $M_B(\beta)$ den Rang $1$, also
\[
  \mathrm{Rg}(\beta) = \mathrm{Rg}(M_B(\beta)) = 1.
\]

\end{document}
